% ICAISD 2026 submission -- compute-aware traffic forecasting benchmark.
%
% FORMAT: see paper/FORMAT_REQUIREMENTS.md. Submission is SINGLE column
% (manuscript mode); ACM production typesets the two-column look after
% acceptance, so do not fight the layout. Minimum 8 pages, NOT a cap.
%
% NUMBERS ARE NOT HAND-TYPED. Table 1 comes from
%   python -c "from greenbench import analysis; \
%              print(analysis.latex_table(analysis.load_results('results')))"
% Anything not yet measured is a \PENDING{} and renders red.
\documentclass[manuscript,screen,review]{acmart}

\usepackage{booktabs}
\usepackage{xcolor}

% Loud placeholder: an unmeasured number must never reach a reviewer quietly.
\newcommand{\PENDING}[1]{\textcolor{red}{\textbf{[PENDING: #1]}}}

\begin{document}

\title{What Does Accuracy Cost? A Compute-Aware Benchmark of
Traffic Forecasting Models on LargeST}

\author{Abhishek Shukla}
\email{shuklabhishek1609@gmail.com}
\affiliation{%
  \institution{Independent Researcher}
  \country{United States}
}
\thanks{Corresponding author.}

\begin{abstract}
Spatio-temporal traffic forecasting is evaluated almost entirely on error
metrics, and the field's leaderboards are therefore silent about what
accuracy costs to obtain. We report a compute-aware benchmark of
six forecasting models on the San Diego subset of LargeST, in which
every model is trained under an identical budget on identical hardware
while training energy, inference energy, FLOPs, parameter count and
inference latency are measured rather than estimated. We find that the
accuracy--cost frontier is extremely uneven: \PENDING{headline comparison,
e.g.\ the most expensive model's MAE advantage over a lightweight
embedding model is 0.41, which is within the measured run-to-run spread,
while costing 45$\times$ the training energy and 83$\times$ the FLOPs}. Two
widely used baselines are strictly dominated, achieving worse error than a
cheaper model at an order of magnitude more energy. We argue that reporting
error alone is insufficient for a field whose models are increasingly
deployed on metropolitan-scale sensor networks, and we release the
instrumented harness so that cost can be reported as routinely as error.
\end{abstract}

\begin{CCSXML}
<ccs2012>
<concept>
<concept_id>10010147.10010257.10010293.10010294</concept_id>
<concept_desc>Computing methodologies~Neural networks</concept_desc>
<concept_significance>500</concept_significance>
</concept>
<concept>
<concept_id>10010520.10010553.10010562</concept_id>
<concept_desc>Computer systems organization~Embedded and cyber-physical systems</concept_desc>
<concept_significance>300</concept_significance>
</concept>
</ccs2012>
\end{CCSXML}

\ccsdesc[500]{Computing methodologies~Neural networks}
\ccsdesc[300]{Computer systems organization~Embedded and cyber-physical systems}

\keywords{traffic forecasting, green AI, energy measurement,
spatio-temporal graph neural networks, benchmarking, sustainable computing}

\maketitle

\section{Introduction}
Traffic forecasting has become one of the standard proving grounds for
spatio-temporal deep learning. A decade of work has produced a dense
sequence of architectures---recurrent, convolutional, graph-based, and
attention-based---each reporting improvements in mean absolute error on a
small number of shared benchmarks. Almost none of that work reports what
the improvement cost.

This omission is not merely an accounting gap. Traffic forecasting is an
applied field: models are proposed for deployment on real sensor networks
operated by transportation agencies, where inference runs continuously
across thousands of sensors and retraining is periodic rather than
one-off. In that setting the relevant question is not only which model is
most accurate but which model is worth its energy, latency and hardware
footprint. A model that reduces MAE by two percent while consuming an
order of magnitude more energy is a different proposition from one that
achieves the same reduction for free, and the published record does not
currently distinguish them.

We take the position that cost should be a reported quantity rather than
an afterthought, and that it should be \emph{measured} rather than
inferred from parameter counts. Parameter count is a poor proxy: it
ignores the sequential structure that makes recurrent models slow, the
dense message passing that makes graph models expensive, and the fact that
models differ in average power draw while running, not only in runtime.

This paper contributes a compute-aware benchmark of six traffic
forecasting models evaluated on the San Diego subset of the LargeST
dataset. Our contributions are:

\begin{itemize}
  \item \textbf{An instrumented evaluation harness} that measures training
        energy, inference energy, FLOPs, parameter count and inference
        latency alongside conventional error metrics, on a single fixed
        GPU, under a training budget held identical across models.
  \item \textbf{A measured accuracy--cost frontier} for the model families
        in common use, identifying which models lie on the Pareto frontier
        and which are strictly dominated.
  \item \textbf{An explicit treatment of measurement validity.} We report
        run-to-run variance, identify which of our own numbers are bounds
        rather than point estimates, and state where a measurement reflects
        the instrument rather than the model.
  \item \textbf{A released harness} so that cost reporting is reproducible
        and cheap for subsequent work to adopt.
\end{itemize}

\section{Related Work}
\subsection{Traffic forecasting benchmarks}
Modern traffic forecasting is organised around a small number of shared
benchmarks. METR-LA and PEMS-BAY, released with the diffusion convolutional
recurrent network of Li et al.~\cite{li2018dcrnn}, established the protocol
that the field still follows: a fixed sensor network, a fixed chronological
split, a twelve-step input window and a twelve-step forecast horizon, scored
by mean absolute error, root mean squared error and mean absolute percentage
error. The architectures that followed were evaluated by dropping them into
that protocol. Spatio-temporal graph convolutional networks~\cite{yu2018stgcn}
replaced recurrence with convolution over the sensor graph; Graph
WaveNet~\cite{wu2019graphwavenet} added a learned adjacency matrix so that the
model was not restricted to the road-network topology. Each reported error
reductions against its predecessors on the same two datasets.

The standardisation this produced is genuine and valuable, and recent work has
tightened it further. BasicTS+~\cite{shao2024basicts} demonstrates that a
substantial part of the apparent progress in multivariate forecasting is an
artefact of inconsistent training pipelines, and supplies a unified pipeline
under which more than forty methods can be compared without that confound.
LargeST~\cite{liu2023largest} attacks a different limitation --- scale ---
observing that METR-LA and PEMS-BAY contain a few hundred sensors and a few
months of data, which is not the regime in which a transportation agency
actually operates. It provides 8,600 California sensors over five years,
together with sensor metadata, and partitions them into regional subsets.

What none of this standardisation covers is cost. BasicTS+ makes the error
protocol fair; LargeST makes it large; neither makes the accounting of what a
result cost to produce a reported quantity. The LargeST authors do report
efficiency observations alongside their baselines, but as a practical remark
about which models are trainable at scale rather than as a measured axis with
a stated instrument. Our contribution sits precisely in that gap: we keep the
error protocol these benchmarks established and add a cost protocol beside it,
with the same seriousness about measurement and hardware control.

\subsection{Efficiency and energy reporting in machine learning}
The argument that efficiency belongs in the reported results rather than in an
appendix is now well developed. Schwartz et al.~\cite{schwartz2020greenai}
name the problem directly, contrasting ``Red AI'' --- accuracy bought with
unreported compute --- against a Green AI in which efficiency is a first-class
evaluation criterion, and propose the floating-point operation count as a
reportable measure. Strubell et al.~\cite{strubell2019energy} give the
empirical counterpart for natural language processing, estimating the energy
and carbon cost of training then-current models and finding costs large enough
to change how a research programme should be judged. Henderson et
al.~\cite{henderson2020systematic} move from estimate to instrument, arguing
that energy should be read from hardware counters rather than inferred, and
providing tooling to do so.

Two refinements of that literature bear directly on our design. The first is
that estimated and measured cost can diverge substantially, because models
differ not only in how long they run but in how much power they draw while
running --- a quantity no analytic FLOP count expresses. The second is
Luccioni et al.'s~\cite{luccioni2024power} observation that for a deployed
system, inference rather than training dominates lifetime cost, since training
happens occasionally and inference happens continuously. Traffic forecasting
is exactly such a deployed setting: a metropolitan sensor network produces
forecasts every few minutes, indefinitely. We therefore measure training
energy and inference energy separately and report both, rather than treating
training cost as a proxy for total cost.

Our contribution is not a new method for measuring energy. It is the
application of this existing discipline to an applied benchmark that has not
adopted it, on hardware held fixed, with the resulting numbers placed beside
the error numbers in the same table.

\subsection{Efficient time-series models}
A parallel line of work questions whether the architectural complexity in this
area is earning its keep. Zeng et al.~\cite{zeng2023transformers} show that
single-layer linear models are competitive with, and often superior to,
elaborate transformer architectures on long-horizon forecasting benchmarks,
and introduce NLinear and DLinear as deliberately minimal baselines. Within
spatio-temporal forecasting specifically, STID~\cite{shao2022stid} makes a
sharper claim: that much of what spatio-temporal graph neural networks
contribute can be recovered by attaching learned spatial and temporal identity
embeddings to a multi-layer perceptron, with no graph convolution at all. The
authors attribute the gap to sample indistinguishability rather than to
relational structure per se.

These results are usually read as accuracy claims --- the simple model matches
the complex one. Our measurements let them be read as cost claims instead,
which is the stronger form of the argument: the question is not merely whether
a lightweight model ties a heavyweight one on error, but what the heavyweight
model charges for the difference, and whether that difference survives the
heavyweight model's own run-to-run variance.

\section{Methodology}
\subsection{Dataset}
We use the San Diego (SD) subset of LargeST~\cite{liu2023largest}, a traffic
sensor dataset derived from California's Performance Measurement System. The
SD subset is chosen deliberately: it is large enough that the cost differences
between architectures are meaningful, while remaining trainable end to end on
a single commodity GPU, which is a precondition for holding hardware fixed
across every model in the comparison.

The subset comprises the 716 sensors in Caltrans District 11. We use the 2019
calendar year, aggregated from the raw five-minute feed into 15-minute bins by
averaging, giving 96 steps per day and 35,040 steps in total; missing readings
are zero-filled. Following the benchmark's protocol we forecast twelve steps
(three hours) from a twelve-step input window, and split the year
chronologically in a 6:2:2 ratio into 21,010 training, 7,003 validation and
7,004 test windows. A chronological rather than random split matters here: the
alternative leaks future traffic conditions into training and inflates every
model's score at once, which would leave the accuracy axis compressed and the
cost comparison harder to interpret.

Each input carries three channels: the traffic flow reading itself, plus
time-of-day and day-of-week indices. This detail is worth stating because it
interacts with one of our results. Of the models we evaluate, NLinear consumes
only the flow channel, since a single linear projection over the input window
has no mechanism for calendar features, whereas STID's entire design premise
is to exploit exactly those identity signals. A reader who does not know this
will read NLinear as a strawman; in fact the comparison between the two
isolates what the calendar embeddings are worth.

Features are standardised using statistics computed on the training split
only, and all reported errors are in the original units of the flow readings.

\subsection{Models}
We evaluate models spanning the cost spectrum rather than only the frontier of
reported accuracy, because the question at issue is the shape of the
trade-off, not the identity of the winner:

\begin{itemize}
  \item \textbf{Historical Last}, a parameter-free baseline that repeats the
        most recent observation across the forecast horizon. It anchors the
        cost axis at zero training energy and supplies the error level that
        any learned model must beat to justify being trained at all.
  \item \textbf{NLinear}~\cite{zeng2023transformers}, a single linear
        projection over the input window, applied after subtracting the last
        observed value and added back before output.
  \item \textbf{STID}~\cite{shao2022stid}, a multi-layer perceptron over
        learned spatial and temporal identity embeddings, with no graph
        convolution.
  \item \textbf{LSTM}, a recurrent baseline that models each sensor's history
        sequentially and shares no information across sensors.
  \item \textbf{STGCN}~\cite{yu2018stgcn}, which alternates temporal
        convolutions with graph convolutions over the road-network adjacency
        matrix.
  \item \textbf{Graph WaveNet}~\cite{wu2019graphwavenet}, which augments
        dilated temporal convolution with a self-adaptive adjacency matrix
        learned from data rather than taken from the road network.
\end{itemize}

Model implementations and their hyperparameters are taken unmodified from the
LargeST reference code, so that no baseline is disadvantaged by our tuning or
lack of it. We had also intended to include STTN, a spatio-temporal
transformer. We omit it for a reason that is itself a datum for this paper: at
its measured per-epoch cost the model would not complete the training budget
within the twelve-hour limit of the compute available to us, and a run
terminated at that limit yields no usable record. A benchmark about the cost
of accuracy is entitled to report that one of its candidates was too expensive
to benchmark, but we prefer to state the omission plainly rather than present
a truncated run beside converged ones.

\subsection{Instrumentation}
\label{sec:instr}
All measurements are taken on a single NVIDIA T4. Holding the device fixed
is essential rather than incidental: energy readings are not comparable
across GPU generations, so a benchmark whose runs are spread over
heterogeneous hardware cannot support cost claims at all. We verify that
every reported run carries the same device identifier before aggregating.

\paragraph{Energy.} Training and inference energy are read from the
device's total-energy counter through NVML, which accumulates joules at the
hardware level. We report training energy as the integral over the whole
training run, and inference energy normalised per one thousand forecasts.

\paragraph{A measurement bug worth reporting.} Our first implementation
measured inference energy over a fixed number of batches. For the cheapest
models this amounted to a few milliseconds of GPU work, far below the
resolution of the energy counter, and the measurement returned exactly
zero joules. A zero inference cost is not a defensible claim, and it
silently distorted the deployment-cost comparison in favour of the
cheapest models. The harness now replays batches until a wall-clock floor
of five seconds has elapsed, with device synchronisation between passes.
We report this because a benchmark's credibility rests on its instrument,
and an unreported instrument failure of this kind is indistinguishable
from a result.

\paragraph{FLOPs, latency and parameters.} FLOPs per sample are counted
analytically and flagged where the count is unreliable. Latency is
measured per forward pass at batch size one after warm-up. Parameter
counts include trainable parameters only.

\subsection{Training protocol}
Every model trains under an identical budget: a maximum of 100 epochs with
early stopping at patience 30, the same optimiser settings, the same seed,
and the same data pipeline. A fixed budget is the coherent choice for a
cost-aware comparison, but it has a consequence that must be stated
plainly: a model that reaches the epoch cap without early stopping has not
converged, so its energy figure is a \emph{lower bound} on its true
cost-to-convergence and its error is pessimistic. We mark such models
explicitly rather than silently reporting their numbers alongside
converged ones.

\section{Results}

Table~\ref{tab:main} reports every model on both axes. Rows are ordered by
training energy, so the cost gradient runs down the table and can be read
against the error columns beside it.

\begin{table*}[t]
\centering
\caption{Accuracy and cost on LargeST-SD (716 sensors, 2019, 15-min). Energy is GPU-only, measured via NVML. Carbon assumes 195~gCO$_2$e/kWh. $\dag$~marks a model that reached the 100-epoch cap without early stopping: its energy is a lower bound, its error pessimistic.}
\label{tab:main}
\begin{tabular}{lrrrrrrrrr}
\toprule
Model & Params & MFLOPs & MAE & RMSE & MAPE (\%) & Train (Wh) & gCO$_2$e & Infer (J/1k) & Latency (ms) \\
\midrule
Historical Last & 0 & -- & 60.79 & 87.40 & 41.88 & 0 & 0 & 0.03 & 0.10 \\
LSTM$^{\dag}$ & 97,932 & 1042.77 & 26.35 & 41.60 & 17.04 & 123.21 & 24.03 & 86.85 & 85.10 \\
STGCN & 507,532 & 1018.76 & 19.78 & 34.18 & 14.01 & 165.05 & 32.19 & 134.11 & 202.63 \\
Graph WaveNet$^{\dag}$ & 311,164 & 11883.93 & 18.30 & 30.61 & 12.39 & 568.00 & 110.78 & 369.25 & 355.04 \\
\bottomrule
\end{tabular}
\end{table*}

\subsection{Accuracy}
\PENDING{Final ordering by MAE/RMSE/MAPE once the lightweight models land;
state the run-to-run spread before the ordering so that differences below it
are not over-read.} Among the models measured so far the ordering is stable
across all three error metrics: Graph WaveNet is most accurate at 18.30 MAE,
followed by STGCN at 19.78, LSTM at 26.35, and Historical Last at 60.79. That
the learned models improve on repeating the last observation by a factor of
roughly three is expected and is reported only to establish that the task is
being learned rather than trivially matched.

The more informative comparison is among the learned models, where the spread
is narrow. Graph WaveNet's advantage over STGCN is 1.47 MAE, or 7.4\% relative.
LSTM is the outlier, 6.57 MAE behind STGCN, and it should be read with the
caveat that it never early-stopped: its error is pessimistic and a longer
budget would improve it somewhat. We return to this in
Section~\ref{sec:threats}. The gap is not, however, of a size that additional
epochs would close.

\subsection{Cost}
The cost axis spans a far wider range than the accuracy axis, and this
asymmetry is the central empirical observation of the paper. Across the models
measured, training energy varies by a factor of 4.61 between LSTM and Graph
WaveNet, while their MAE differs by 8.05 --- but between STGCN and Graph
WaveNet, where the error difference has narrowed to 1.47 MAE, the energy
difference is still 3.44$\times$. The pattern is that error differences
compress as models improve while cost differences do not. \PENDING{Restate
these ratios over the full model set, including the lightweight models, which
extend the low end of the cost axis by roughly an order of magnitude.}

The deployment-side axes are wider still. Inference energy per thousand
forecasts ranges from 0.03~J for the parameter-free baseline to 369~J for
Graph WaveNet, a factor of 14,563, and single-sample latency from 0.10~ms to
355~ms, a factor of 3,550. For a network of several hundred sensors producing
forecasts continuously, these are the figures that determine the operating
cost of a deployment, and they are not derivable from the error columns.

\paragraph{Parameter count is not a cost proxy.} The clearest single
demonstration in our measurements is the comparison between STGCN and Graph
WaveNet. STGCN has 507,532 trainable parameters against Graph WaveNet's
311,164 --- 1.63$\times$ as many --- and yet requires 11.7$\times$
\emph{fewer} floating-point operations per sample, trains on 3.44$\times$ less
energy, and answers 1.75$\times$ faster. A reader ranking these two models by
parameter count would order them exactly backwards on every cost axis we
measure. This is not a subtle effect at the margin; it is a reversal, and it
is why we regard a measured cost protocol as necessary rather than merely
desirable.

\paragraph{Power draw, not merely runtime.} Cheaper models also draw less
power \emph{while running}, so their advantage compounds beyond their shorter
runtime. \PENDING{Report measured average power per model, computed as
train\_joules / train\_seconds, once the lightweight runs land --- the pilot
measurements showed roughly 31~W for NLinear against 65~W for LSTM on a device
capped at 70~W.} An analytic FLOP count cannot capture this, which is a
further argument for measurement over estimation.

\subsection{The accuracy--cost frontier}
\PENDING{Identify the Pareto frontier and its knee over the full model set,
and name the strictly dominated models --- those achieving worse error at
greater cost. A strictly dominated widely-used baseline is a result in itself,
and on the partial results it is already visible that the frontier is not
monotone in architectural sophistication.}

\subsection{Forecast horizon}
Whether the ranking is stable across the forecast horizon matters for
deployment: a model that wins only at short horizons is a different
recommendation from one that wins throughout. It is also where the
parameter-free baseline's behaviour is most revealing. Historical Last
degrades from 17.72 MAE at the first step to 101.74 at the twelfth --- a
factor of 5.74 --- because repeating the last observation carries no
information about how traffic evolves. The learned models degrade far more
gently: Graph WaveNet from 12.51 to 22.55, a factor of 1.80, and STGCN from
15.31 to 23.42, a factor of 1.53.

This produces a crossover worth noting. At the first horizon step STGCN's
15.31 MAE is worse than Graph WaveNet's 12.51 by 2.80, but by the twelfth step
the gap has narrowed to 0.87. The most expensive model's advantage is
concentrated at short horizons and erodes as the forecast lengthens, which
weakens the case for its cost precisely in the long-horizon regime where
forecasting is most useful operationally. \PENDING{Extend this analysis to the
lightweight models and state whether the frontier ordering is horizon-stable.}

\section{Discussion}
Read together, the two axes of Table~\ref{tab:main} tell a different story
than either would alone. On the error axis the learned models are separated by
single units of MAE, and the separations shrink as the models get better. On
the cost axis they are separated by multiplicative factors that do not shrink.
A field that reports only the first axis sees a sequence of incremental
improvements; a field that reports both sees those improvements being bought
at a steeply rising price.

This has a direct consequence for how architectural claims should be
evaluated. When the accuracy difference between two models is comparable to
the run-to-run variance of the training procedure that produced them, while
the cost difference is an order of magnitude, the burden of proof shifts. A
new architecture claiming a fractional MAE improvement should be expected to
report what the improvement cost, and to demonstrate that the improvement
exceeds the variance of its own training procedure --- a demonstration that
requires more than one seed and is, at present, rarely offered. Our own
measurements make this concrete: we observed a spread of approximately
$\pm0.3$ MAE between two runs of the same model under the same seed, which is
a substantial fraction of the differences separating the leading models in
Table~\ref{tab:main}.

The parameter-count reversal between STGCN and Graph WaveNet sharpens the
point. Parameter count is the cost proxy most often reported when any is
reported at all, and on our measurements it orders those two models backwards
on FLOPs, on training energy and on latency simultaneously. A proxy that can
invert the ranking it is standing in for is not a weak measurement; it is a
misleading one. Nothing short of instrumenting the run recovers the true
ordering.

None of this argues that expensive models are never worth their cost. It
argues that the question is currently unanswerable from the published record,
because one side of the ledger is missing. Where the cost is worth paying ---
a safety-critical application, or a short-horizon forecast where the accuracy
gap is widest --- a practitioner should pay it, and our measurements let that
judgement be made rather than guessed.
\PENDING{Close with the frontier claim once the lightweight models land:
state the knee and what it implies for a practitioner choosing a model for a
metropolitan deployment.}

\paragraph{What it would take to adopt this.} The instrumentation reported
here is not elaborate. Energy comes from a counter the driver already
maintains; FLOPs and parameters are counted analytically; latency requires a
warm-up loop and a device synchronisation. The measurement bug described in
Section~\ref{sec:instr} was the only genuine difficulty we encountered, and it
is a solved problem now that it is documented. The cost of adding a cost
protocol to an existing benchmark is a few hours of engineering, against which
the cost of not having one is a literature whose central trade-off is
invisible.

\section{Threats to Validity}
\label{sec:threats}
\paragraph{Single seed.} Results are reported for one seed per model. We
measured the run-to-run spread directly at approximately $\pm0.3$ MAE, by
comparing two runs of the same model under identical seed and configuration,
and we treat accuracy differences smaller than that spread as unresolved
rather than as rankings. Our training is not bit-reproducible: the harness
does not force deterministic cuDNN kernels, since doing so changes the
performance characteristics we are trying to measure. This is a real
limitation of the accuracy axis. It does not materially affect the cost axis,
where the differences we report are multiplicative and orders of magnitude
larger than any plausible run-to-run variation in energy.

\paragraph{Models that hit the epoch cap.} LSTM and Graph WaveNet both reached
the 100-epoch cap without early stopping; Graph WaveNet's validation error was
still improving at epoch 98. For these two models the reported training energy
is a lower bound on the true cost to convergence, and the reported error is
pessimistic. They are marked in Table~\ref{tab:main} rather than mixed
silently with converged runs. This cuts against our own argument rather than
for it: Graph WaveNet's cost is understated, not overstated, so correcting it
would widen the cost gap we report rather than narrow it.

\paragraph{The parameter-free baseline's inference energy.} Historical Last
has no computation to perform, so its measured inference energy of 0.03~J per
thousand forecasts is dominated by idle device draw during the measurement
window. This figure is the noise floor of the instrument, not a property of
the model, and the 14,563$\times$ inference-energy span quoted against it
should be read as the dynamic range of the measurement rather than as a clean
model-to-model ratio.

\paragraph{Energy scoping.} We measure GPU energy only. Host CPU, system
memory, storage, networking, cooling and data-centre overhead are excluded, so
every energy and carbon figure reported here is a lower bound on the true
footprint of the run. We prefer this scoping to a whole-system estimate
multiplied by an assumed power usage effectiveness, because the GPU counter is
measured while the multiplier would be assumed, and mixing the two would
obscure which part of the number is evidence.

\paragraph{Single dataset and single device.} Results characterise one subset
of one dataset on one GPU, an NVIDIA T4. Energy readings are not comparable
across GPU generations, and we verified that every reported run carries the
same device identifier before aggregating. The cost \emph{ratios} we report
are more portable than the absolute joule figures, but neither should be
extrapolated to other hardware without re-measurement.

\paragraph{Carbon figures are derived, not measured.} Emissions are computed
from measured energy using a fixed grid intensity of 195~gCO$_2$e/kWh, derived
from the United States Environmental Protection Agency's eGRID subregion CAMX
(WECC California) total output rate of 430.0~lb~CO$_2$e/MWh in eGRID2023
Revision 2~\cite{epa2025egrid}. We use the California figure because LargeST
is California data and a deployment on this sensor network would draw on this
grid. The constant carries the usual caveats: it is an annual average over a
grid whose real carbon intensity varies by hour and by season, so the carbon
column should be read as a scaled version of the energy column rather than as
an independent measurement.

\section{Conclusion}
We have reported a compute-aware benchmark of six traffic forecasting models
on the San Diego subset of LargeST, in which every model trains under an
identical budget on identical hardware while training energy, inference
energy, FLOPs, parameters and latency are measured rather than estimated.
Accuracy and cost are not proportional on this benchmark. Error differences
between the leading models compress into single units of MAE, comparable to
the run-to-run variance of the training procedures that produced them, while
cost differences remain multiplicative across every axis we measure --- up to
14,563$\times$ in inference energy and 3,550$\times$ in latency. Parameter
count, the proxy most often reported in place of a measurement, inverts the
true cost ordering of two of our models on three separate axes.
\PENDING{Add the frontier and dominance claims once the lightweight models
land.} We release the instrumented harness so that reporting cost becomes as
routine as reporting error, and so that the next architecture proposed for
this benchmark can be asked not only whether it is more accurate, but what its
accuracy costs.

\paragraph{Reproducibility.} The harness, the run configurations and the raw
per-run records underlying Table~\ref{tab:main} are available at
\url{https://github.com/shuklabhishek1609-blip/ICAISD}.

\begin{acks}
Use of AI: the author used a large language model as a coding and drafting
assistant during the development of the measurement harness and the
preparation of this manuscript. All experimental design decisions were made
by the author; every reported measurement was produced by executing the
described harness on the stated hardware; and no result in this paper was
generated, estimated or extrapolated by a language model.
\end{acks}

\bibliographystyle{ACM-Reference-Format}
\bibliography{refs}

\end{document}
